\section{Отношения}

\begin{definition}[Отношение]
    Любое подмножество $R$ декартова произведения $k$ множеств, для $k > 1, k < \infty$, называется \emph{$k$-местным отношением}.
\end{definition}

Если есть $k$-местное отношение $R$ и $(x_0,\ldots, x_{k-1}) \in R$, то часто говорят, что $x_i$ \emph{находятся в отношении} $R$.

Если $k=2$ и $A_0 = A_1 = A$, то говорят о \emph{бинарных отношениях} над $A$.

\begin{definition}[Рефлексивное отношение]
    Бинарное отношение $R$ над $A$ является \emph{рефлексивным}, если для любого $a \in A$ $(a,a) \in R$.
\end{definition}


\begin{definition}[Транзитивное отношение]
    Бинарное отношение $R$ над $A$ является \emph{транзитивным}, если для любых $a_0 \in A, a_1 \in A, a_2 \in A$, если $(a_0,a_1) \in R$ и  $(a_1,a_2) \in R$, то $(a_0,a_2) \in R$.
\end{definition}

\begin{definition}[Транзитивное замыкание отношения]
    Пусть $R$ --- бинарное отношение над $A$.
    Тогда транзитивным замыканием $R$ является наименьшее по мощности транзитивное отношение $R_t$ над $A$, такое, что\sidenote{Помним, что отношение --- это тоже множество.} $R \subseteq R_t$.
\end{definition}

\begin{definition}[Рефлексивно-транзитивное замыкание отношения]
    Пусть $R$ --- бинарное отношение над $A$.
    Тогда рефлексивно-транзитивным замыканием $R$ является наименьшее по мощности транзитивное и рефлексивное отношение $R_{rt}$ над $A$, такое, что $R \subseteq R_{rt}$.
\end{definition}

\begin{definition}[Композиция отношений]
    Пусть $R_1 \subseteq A \times B$, $R_2 \subseteq B \times C$ --- бинарные отношения.
    Тогда композицией отношений $R_1$ и $R_2$ будем называть бинарное отношение $R_3 = \{ (a,c) \mid (a,b) \in R_1 , (b,c) \in R_2\}$.
    Для обозначения композиции будем использовать следующую запись: $R_3 = R_1 \circ R_2$.
\end{definition}

Пусть есть бинарное отношение $R \subseteq A \times B$, где $A$ и $B$ --- конечные множества.
Дополнительно будем считать, что оба этих множества проиндексированы.
Тогда $R$ можно задать с помощью булевой матрицы $M_R$ размера $|A|\times|B|$ следующим образом: $M[i,j] = 1 \iff (A[i],B[j] \in R)$.

\begin{theorem}
    Пусть $R_1 \subseteq A \times B$, $R_2 \subseteq B \times C$ --- бинарные отношения, а множества $A$, $B$ и $C$ конечны.
    $M_{R_1}$ и $M_{R_2}$ --- булевы матрицы, задающие соответствующие отношения.
    Тогда $M_{R_3} = M_{R_1} \times M_{R_2}$ --- матрица, задающая отношение $R_3 = R_1 \circ R_2$.
\end{theorem}

Доказательство данной теоремы выполняется на основе определения умножения матриц\sidenote{Оставим его как упражнение для читателя.}.
